\documentclass[a4paper,11pt,oneside]{book}
% preamble_build.tex: copy of Maximus's preamble.tex for the DE v3 build.
% The blocks marked [BUILD COPY CHANGE] are the earlier build fixes.
% The blocks marked [V3 ADDITION] are new in this version (example box,
% spacing). Everything else is unchanged from preamble.tex.

% For margins
\usepackage[a4paper, margin = 20mm]{geometry}

\usepackage[Lenny]{fncychap}
% Rejne, Glenn, Lenny

\usepackage{etoc}

% [V3 ADDITION] cancel: used to show the dx's "cancelling" in Section 11.3
\usepackage{cancel}
\renewcommand{\CancelColor}{\color{red}}

% For enumeration & math symbols
\usepackage{enumerate}
\usepackage{amsmath,amsfonts,amssymb,amsthm,graphicx,lipsum}

% For diagrams, plots and boxes
\usepackage{pgfplots}
\pgfplotsset{compat=1.18}
\usepackage{tikz}
\usepackage[most]{tcolorbox}
\tcbuselibrary{skins}
\usepackage{xcolor}
\pgfplotsset{
    standard/.style={clip = false,
    axis line style = thick,
    trig format=rad,
    axis lines = middle,
    every axis x label/.style={at={(current axis.right of origin)},anchor= west},
    every axis y label/.style={at={(current axis.above origin)},anchor=south}
    }
}

% For GC font
\usepackage{fontspec}
% [BUILD COPY CHANGE 1] The TI-84 key font is loaded only if the file exists
% (checked in ./Fonts/... and ../Fonts/...). Otherwise \ti{} falls back to a
% small boxed sans-serif key. Known letters of the TI font are mapped to their
% key names in the fallback: o = [y=], z = [mode], f = [L3].
\newcommand{\tikey}[1]{\tcbox[on line, enhanced, boxsep=0pt, left=2pt, right=2pt,
  top=1pt, bottom=1pt, boxrule=0.5pt, arc=1.5pt, colback=white,
  colframe=black!70]{\sffamily\footnotesize\strut #1}}
\ExplSyntaxOn
\NewDocumentCommand{\tifallback}{m}{
  \str_case:nnF {#1}
    { {o}{\tikey{y=}} {z}{\tikey{mode}} {f}{\tikey{L3}} }
    { \tikey{#1} }
}
\ExplSyntaxOff
\def\tifontcase{0}
\IfFileExists{Fonts/TI-84 Plus CE Key Fonts/TI84PlusCEKeys.ttf}
  {\def\tifontcase{1}}{}
\IfFileExists{../Fonts/TI-84 Plus CE Key Fonts/TI84PlusCEKeys.ttf}
  {\def\tifontcase{2}}{}
\ifcase\tifontcase\relax
  % 0: font not found, use boxed fallback keys
  \newcommand{\ti}[1]{\tifallback{#1}}
\or
  % 1: original location (as in preamble.tex)
  \newfontfamily\tifont[
    Path = {Fonts/TI-84 Plus CE Key Fonts/},
    UprightFont = {TI84PlusCEKeys.ttf}
  ]{}
  \newcommand{\ti}[1]{{\tifont #1}}
\or
  % 2: one level up (when compiling inside draft-ch1/)
  \newfontfamily\tifont[
    Path = {../Fonts/TI-84 Plus CE Key Fonts/},
    UprightFont = {TI84PlusCEKeys.ttf}
  ]{}
  \newcommand{\ti}[1]{{\tifont #1}}
\fi

% [FONT SETUP copied from h2-math-notes/aoi/notes/preamble_build.tex (DI notes)]
% [AoI items 17/18/37] GC DISPLAY FONT. Every GC key, command, menu item and
% screen line goes through \gcfont, so the font can be swapped HERE in one place.
%   \gc{fnInt(}      command / menu item / stored variable in the display font
%   \gckey{2nd}      a calculator key (display font in a small key outline)
%   \gcscreen{...}   a calculator screen (display font on a grey-green LCD panel)
% The font file lives in ../Fonts (shared by notes, slides and worksheet).
\def\gcfontfile{VT323-Regular.ttf}
\IfFileExists{Fonts/\gcfontfile}{%
  \newfontfamily\gcfont[Path=Fonts/, Scale=1.32]{\gcfontfile}}{%
  \newcommand\gcfont{\ttfamily}}
\newcommand{\gc}[1]{\mbox{\gcfont #1}}
\newcommand{\gckey}[1]{\tcbox[on line, enhanced, boxsep=0pt, left=2pt, right=2pt,
  top=0.6pt, bottom=0.2pt, boxrule=0.5pt, arc=1.5pt, colback=black!4,
  colframe=black!65]{\gcfont\strut #1}}
\definecolor{lcd}{RGB}{214,222,204}
\newcommand{\gcscreen}[1]{\tcbox[on line, enhanced, boxsep=1pt, left=4pt, right=4pt,
  top=2pt, bottom=2pt, boxrule=1.2pt, arc=1pt, colback=lcd, colframe=black!70]{%
  \gcfont\begin{tabular}{@{}l@{}}#1\end{tabular}}}
% integral template as shown on the TI-84 Plus CE (MathPrint): \gcint{0}{2}{e^(-X²)}{X}
\newcommand{\gcint}[4]{{\gcfont\raisebox{-0.35ex}{\Large ∫}\hspace{-0.15em}%
  \raisebox{-0.2ex}{\begin{tabular}{@{}l@{}}\footnotesize #2\\[-1.3ex]\footnotesize #1\end{tabular}}%
  \hspace{0.1em}(#3)d#4}}
% old key macro kept as an alias so nothing breaks
\renewcommand{\ti}[1]{\gckey{#1}}

% SI units
\usepackage{siunitx}

% setspace
%\usepackage{setspace}
%\onehalfspacing

% For hyperlinks within and out of document
\usepackage{hyperref}
\hypersetup{
    colorlinks = true,
    linkcolor = black,
    urlcolor = blue,
}   

% For captions for pics
\usepackage{capt-of}


% tcolorbox custom environments
% Definition ========================================================
\newtcolorbox[auto counter,number within=section]{definition}[2][drop shadow]{%
enhanced, before skip=5mm, after skip=6mm, % [V3 ADDITION] spacing
colback=blue!5!white,colframe=blue!50!black,fonttitle=\bfseries,
title= Definition~\thetcbcounter. #2,#1}

% Theorem ==========================================================
\newtcolorbox[auto counter,number within=section]{theorem}[2][drop shadow]{%
enhanced, before skip=5mm, after skip=6mm, % [V3 ADDITION] spacing
colback=black!5!white,colframe=black,fonttitle=\bfseries,
title= Theorem~\thetcbcounter. #2,#1}

% Caution box ======================================================
\newtcolorbox{caution}[1][]{enhanced,
before skip=5mm,after skip=6mm, % [V3 ADDITION] was 2mm/3mm
boxrule=1pt,left=11mm,right=2mm,top=3mm,bottom=3mm,
colback=yellow!50,
colframe=yellow!20!black,
sharp corners,rounded corners=southeast,arc is angular,arc=3mm,
underlay={%
\path[fill=tcbcolback!80!black] ([yshift=3mm]interior.south east)--++(-0.4,-0.1)--++(0.1,-0.2);
\path[draw=tcbcolframe,shorten <=-0.05mm,shorten >=-0.05mm] ([yshift=3mm]interior.south east)--++(-0.4,-0.1)--++(0.1,-0.2);
\path[fill=yellow!50!black,draw=none] (interior.south west) rectangle node[white]{\Huge\bfseries !} ([xshift=9mm]interior.north west);
},
drop fuzzy shadow,#1}

% Question box ======================================================
\newtcolorbox{question}[1][]{enhanced,
before skip=5mm,after skip=6mm, % [V3 ADDITION] was 2mm/3mm
boxrule=1pt,left=11mm,right=2mm,top=3mm,bottom=3mm,
colback=blue!20!white,
colframe=blue!20!black,
sharp corners,rounded corners=southeast,arc is angular,arc=3mm,
underlay={%
\path[fill=tcbcolback!80!black] ([yshift=3mm]interior.south east)--++(-0.4,-0.1)--++(0.1,-0.2);
\path[draw=tcbcolframe,shorten <=-0.05mm,shorten >=-0.05mm] ([yshift=3mm]interior.south east)--++(-0.4,-0.1)--++(0.1,-0.2);
\path[fill=blue!50!black,draw=none] (interior.south west) rectangle node[white]{\Huge\bfseries ?} ([xshift=9mm]interior.north west);
},
drop fuzzy shadow,#1}

% Information box ======================================================
\newtcolorbox{info}[1][]{enhanced,
before skip=5mm,after skip=6mm, % [V3 ADDITION] was 2mm/3mm
boxrule=1pt,left=11mm,right=2mm,top=3mm,bottom=3mm,
colback=cyan!20!white,
colframe=cyan!20!black,
sharp corners,rounded corners=southeast,arc is angular,arc=3mm,
underlay={%
\path[fill=tcbcolback!80!black] ([yshift=3mm]interior.south east)--++(-0.4,-0.1)--++(0.1,-0.2);
\path[draw=tcbcolframe,shorten <=-0.05mm,shorten >=-0.05mm] ([yshift=3mm]interior.south east)--++(-0.4,-0.1)--++(0.1,-0.2);
\path[fill=cyan!50!black,draw=none] (interior.south west) rectangle node[white]{\Huge\bfseries i} ([xshift=9mm]interior.north west);
},
drop fuzzy shadow,#1}

% Mistake box ======================================================
\newtcolorbox{mistake}[1][]{enhanced,
before skip=5mm,after skip=6mm, % [V3 ADDITION] was 2mm/3mm
boxrule=1pt,left=11mm,right=2mm,top=3mm,bottom=3mm,
colback=red!20!white,
colframe=red!20!black,
sharp corners,rounded corners=southeast,arc is angular,arc=3mm,
underlay={%
\path[fill=tcbcolback!80!black] ([yshift=3mm]interior.south east)--++(-0.4,-0.1)--++(0.1,-0.2);
\path[draw=tcbcolframe,shorten <=-0.05mm,shorten >=-0.05mm] ([yshift=3mm]interior.south east)--++(-0.4,-0.1)--++(0.1,-0.2);
\path[fill=red!50!black,draw=none] (interior.south west) rectangle node[white]{\Huge\bfseries $\times$} ([xshift=9mm]interior.north west);
},
drop fuzzy shadow,#1}



% [V3 ADDITION] Worked example box ==================================
% Own counter, numbered within the section (Example 11.2.1, 11.2.2, ...).
% For per-subsection numbering (11.2.1.1) change `section' to `subsection'. Blue `question' boxes are kept for thought-provoking questions
% only; worked examples always use this box. Usage:
%   \begin{example} ...question text... \end{example}
%   \solution{...}
\newtcolorbox[auto counter,number within=section]{example}[1][]{%
enhanced, breakable=false, before skip=5mm, after skip=4mm,
colback=teal!5!white, colframe=teal!55!black, boxrule=0.8pt,
left=3mm, right=3mm, top=2mm, bottom=2mm,
sharp corners, rounded corners=northwest, arc=3mm,
fonttitle=\bfseries, coltitle=white,
title=Example~\thetcbcounter, attach boxed title to top left={yshift=-2mm, xshift=3mm},
boxed title style={colback=teal!55!black, sharp corners, rounded corners=northwest, arc=2mm},
drop fuzzy shadow, #1}

% [V3 ADDITION] Spacing between sections, subsections, parts, paragraphs =
% (fncychap's chapter heading is untouched; only \section and below.)
\makeatletter
\renewcommand\section{\@startsection{section}{1}{\z@}%
  {-2ex \@plus -.5ex \@minus -.2ex}{3.2ex \@plus .3ex}%
  {\normalfont\LARGE\bfseries\color{blue!50!black}}}
\renewcommand\subsection{\@startsection{subsection}{2}{\z@}%
  {-1.5ex \@plus -.5ex \@minus -.2ex}{2.6ex \@plus .3ex}%
  {\normalfont\Large\bfseries}}
\renewcommand\subsubsection{\@startsection{subsubsection}{3}{\z@}%
  {-4ex \@plus -1ex \@minus -.5ex}{1.6ex \@plus .2ex}%
  {\normalfont\large\bfseries\color{teal!50!black}}}
\makeatother
\setlength{\parindent}{0pt}
\setlength{\parskip}{0.9ex plus 0.3ex minus 0.2ex}

% Headers and Footers
\usepackage{fancyhdr,lastpage}
% [V3 R2] grey left-margin line on the out-of-syllabus pages (item 13)
\usepackage{eso-pic}
\newcommand{\oosmargin}{\AddToShipoutPictureBG{\put(12mm,0){\color{gray!75}\rule{1.6pt}{\paperheight}}}}
\pagestyle{fancy}
\lhead{Maximus Pung Jun}
% [BUILD COPY CHANGE 2] Header bug fix: \chaptername is the word "Chapter",
% so the original header printed "Chapter 1: Chapter". \chaptermark now stores
% the bare chapter title in \leftmark, and the header uses \leftmark.
\renewcommand{\chaptermark}[1]{\markboth{#1}{}}
\rhead{Chapter \thechapter: \leftmark}
% \cfoot{Page \thepage\  of \pageref{LastPage}} % buggy, does not appear on chapter pages.

% [V3 R2] tick and cross markers for mistake boxes (G3)
\newcommand{\wrong}{\textcolor{red!70!black}{\ensuremath{\boldsymbol{\times}}}}
\newcommand{\cmark}{\textcolor{green!50!black}{\ensuremath{\checkmark}}}

\newcommand{\xmark}{\wrong}   % alias: red cross
% [V3 R2] out-of-syllabus helpers (DE notes item 13)
%   \oosmargin    : grey line in the left margin on every following page
%   \oosmarginoff : stop the margin line
%   \oosnote      : grey banner "Out of syllabus. For your interest only."
%   outofsyllabus : environment = new page + margin line + banner; ends with
%                   \clearpage and switches the margin line off again
\newcommand{\oosmarginoff}{\ClearShipoutPictureBG}
\newcommand{\oosnote}[1][Nothing in this section is examined in H2 Mathematics.]{\par\noindent\colorbox{gray!15}{\parbox{\dimexpr\textwidth-2\fboxsep}{\textbf{Out of syllabus. For your interest only.} #1}}\par}
\newenvironment{outofsyllabus}{\clearpage\oosmargin}{\clearpage\oosmarginoff}

% Notation as in Maximus's notes: italic d, display style everywhere (G1).
% Fractions inside exponents stay small, so never use these macros inside ^{}.
\newcommand{\dydx}{\dfrac{dy}{dx}}
\newcommand{\dvdx}{\dfrac{dv}{dx}}
\newcommand{\dudx}{\dfrac{du}{dx}}
\newcommand{\ddydx}{\dfrac{d^{2}y}{dx^{2}}}
\newcommand{\ddt}[1]{\dfrac{d#1}{dt}}
% Small syllabus flag (no longer used in the DE notes, kept for other chapters)
\newcommand{\syllabusflag}[1]{\par\noindent{\footnotesize\color{red!70!black}\textbf{Syllabus note.} #1}\par}


% Teacher's copy toggle ================================================
% [BUILD COPY CHANGE 3] \ifx compares the whole word (the original
% \if\conform test was always false). DE_notes_student.tex predefines
% \conform as HideSolutions.
% A student copy defines \conform as HideSolutions BEFORE loading this preamble.
\providecommand\conform{ShowSolutions}
\newcommand\ShowSolutionsKey{ShowSolutions}
% [V3 ADDITION] In the student copy the hidden solution leaves stretchable
% blank space (\vfill), so the free space on each page is shared out as
% working space below the examples instead of collecting at the bottom.
\newcommand\solution[1]{\ifx\conform\ShowSolutionsKey
  \par\noindent\textbf{Solution.} #1\par\medskip\else\par\vspace{\workspace}\fi}
% [V3 R2] minimum working space left under each example in the student copy
\newlength\workspace\setlength\workspace{3cm plus 1fill}


\newcommand{\NN}{\mathbb{N}}
\newcommand{\ZZ}{\mathbb{Z}}
\newcommand{\QQ}{\mathbb{Q}}
\newcommand{\RR}{\mathbb{R}}
\newcommand{\CC}{\mathbb{C}}


% =====================================================================
\usepackage{pifont}
% preamble_local.tex: extra packages and styles for the Vectors chapter
% ONLY. Loaded after preamble_build.tex (a verbatim copy of the shared
% book preamble from the PnC / DE build, so that it can be re-synced by
% copying it over again).
% =====================================================================
\usepackage{array, booktabs}
\usetikzlibrary{arrows.meta, calc, positioning, 3d, decorations.pathreplacing, angles, quotes}
% =====================================================================
% vec_common.tex: shared by the notes (via notes/preamble_local.tex) and
% the slides (via slides/slides_local.tex) for the Vectors chapter.
%
%   CHAPTER NUMBER: change the 13 below and rebuild. The notes set the
%   chapter counter from it, and every slide frame title, Example/Quiz
%   number and box number on the slides uses \chapnum.
% =====================================================================
\def\chapnum{13}

% ---- answers -------------------------------------------------------------
% \ans{key}{value} just prints the value. checks/verify.py reads every
% \ans in the notes and slides and checks it with sympy.
\newcommand{\ans}[2]{#2}

% Column vector, the form used in MF27 and in the examples.
% pmatrix cells are text style, so \dfrac alone still looks small and the
% rows collide. Force display style in every cell and open the rows up.
% Each cell gets a strut as tall as its own content plus \colpad above and
% below, so rows of \dfrac never touch, while integer rows stay compact.
\newlength{\colpad}\setlength{\colpad}{0.55ex}
\newcommand{\colcell}[1]{%
  \setbox0=\hbox{$\displaystyle #1$}%
  \rule[-\dimexpr\dp0+\colpad\relax]{0pt}{\dimexpr\ht0+\dp0+2\colpad\relax}%
  \box0}
\newcommand{\col}[3]{%
  \begingroup
  \renewcommand{\arraystretch}{1.35}%
  \begin{pmatrix}
    \colcell{#1}\\
    \colcell{#2}\\
    \colcell{#3}
  \end{pmatrix}%
  \endgroup
}

% \place{Name}{x}{y}{z} defines a TikZ coordinate. The parentheses are
% inside the macro, so TikZ sees them only after expansion, and
% checks/verify.py can read the drawn coordinates back out.
\newcommand{\place}[4]{\coordinate (#1) at (#2,#3,#4);}

% ---- 3D pictures ----------------------------------------------------------
% Oblique axes: x to the right and slightly down, y out of the page,
% z straight up. Use scale to fit the frame.
\tikzset{
  threeD/.style={
    x={(0.78cm,-0.20cm)},
    y={(0.36cm,0.26cm)},
    z={(0cm,0.78cm)},
    line cap=round, line join=round,
    >=Stealth
  },
  varr/.style={->, thick},
  pt/.style={circle, fill=black, inner sep=1.15pt},
  plane/.style={fill=blue!12, draw=blue!45!black, thick, fill opacity=0.85},
  hid/.style={dashed, black!55}
}
% vecfig.tex: shared 3D diagram style (tikz-3dplot) for Vectors I / II notes and slides.
\usepackage{tikz-3dplot}
\tdplotsetmaincoords{62}{55}
\definecolor{vA}{RGB}{31,90,180}   % first vector / line 1
\definecolor{vB}{RGB}{210,90,20}   % second vector / line 2
\definecolor{vC}{RGB}{0,128,110}   % result / third
\definecolor{vN}{RGB}{140,40,150}  % normals
\definecolor{vP}{RGB}{200,215,240} % plane 1
\definecolor{vQ}{RGB}{245,215,190} % plane 2
\definecolor{vR}{RGB}{200,235,215} % plane 3
\tikzset{
  v3/.style={tdplot_main_coords, line cap=round, line join=round, >=Stealth, scale=#1},
  v3/.default=1,
  ax/.style={->, black!50, thin},
  va/.style={->, very thick, vA},
  vb/.style={->, very thick, vB},
  vc/.style={->, very thick, vC},
  vn/.style={->, very thick, vN},
  ln/.style={very thick},
  dsh/.style={dashed, black!55},
  pp/.style={fill=vP, fill opacity=0.55, draw=vA!70!black, thin},
  pq/.style={fill=vQ, fill opacity=0.55, draw=vB!70!black, thin},
  pr/.style={fill=vR, fill opacity=0.55, draw=vC!70!black, thin},
  dot/.style={circle, fill=black, inner sep=1.2pt},
}
% \vaxes{xlen}{ylen}{zlen}
\newcommand{\vaxes}[3]{%
  \draw[ax] (0,0,0) -- (#1,0,0) node[anchor=west]{$x$};
  \draw[ax] (0,0,0) -- (0,#2,0) node[anchor=south west]{$y$};
  \draw[ax] (0,0,0) -- (0,0,#3) node[anchor=south]{$z$};
  \node[anchor=east] at (0,0,0) {\scriptsize $O$};}
% \rightmark{corner}{along1}{along2} small right-angle square
\newcommand{\rtang}[3]{\draw[thin] ($(#1)!0.25cm!(#2)$) -- ($($(#1)!0.25cm!(#2)$)+($(#1)!0.25cm!(#3)$)-(#1)$) -- ($(#1)!0.25cm!(#3)$);}
\newenvironment{vis}{\par\medskip\begin{center}}{\end{center}\medskip}
% \vang{vertex}{p1}{p2}{label}: small angle arc drawn on the page
\newcommand{\vang}[4]{%
  \coordinate (vg1) at ($(#1)!0.55cm!(#2)$); \coordinate (vg2) at ($(#1)!0.55cm!(#3)$);
  \coordinate (vgm) at ($(vg1)!0.5!(vg2)$);
  \draw[vC, thick] (vg1) .. controls ($(#1)!0.72cm!(vgm)$) .. (vg2);
  \node[vC] at ($(#1)!0.95cm!(vgm)$) {#4};}


% A little air between the parts of an example, and between lines of a list.
\setlength{\itemsep}{0.85ex}
\setlength{\parsep}{0.35ex}

% Student copy: minimum working space under each example. 2.2cm (instead of
% the book's 3cm) is meant to keep the student copy on the same page breaks
% as the teacher copy.
\setlength\workspace{2.2cm plus 1fill}

% =====================================================================
% prelim_fonts.tex: fonts of the real 9758 prelim papers (Word + MathType; pdffonts on
% ACJC 2026 P1/P2, ACJC 2025 P1, JPJC 2026, RVHS 2026, CJC/DHS/TJC/NYJC/HCI/EJC 2025):
%   text Times New Roman; cover/footer Arial; RI page footer Calibri (here Carlito);
%   maths = MathType: Times New Roman letters/digits (italic variables; upright digits,
%   function names, d, e, f), SymbolMT for the integral, sigma, pi and = + < >,
%   scripts 7 pt on 12 pt (about 58 %), tight operator spacing.
% Candidate C (teacher-approved 13 Oct 2026): the real Times New Roman and Arial
% (ttf-mscorefonts-installer) + URW Standard Symbols PS (fonts-urw-base35, a SymbolMT
% clone with the same Symbol encoding) for the integral, sigma, pi, = + < >;
% TeX Gyre Termes Math only for what MathType draws itself (radicals, stretchy brackets,
% fraction bars, arrows). House style: upright d (\mathrm{d}x), upright e and f.
% Loaded once by di_style / de_core (HW sheets) and prelim_style (Tests, solutions).
% =====================================================================
\ifdefined\prelimfontsloaded\endinput\fi
\def\prelimfontsloaded{1}
\edef\pfrestoreat{\catcode`\noexpand\@=\the\catcode`\@\relax}\makeatletter
% MathType script sizes: 7 pt superscripts on 12 pt text (58 %), 5 pt second level
\DeclareMathSizes{12}{12}{7}{5}
\DeclareMathSizes{11}{11}{6.5}{5}
\DeclareMathSizes{10.95}{10.95}{6.5}{5}
\DeclareMathSizes{10}{10}{6}{5}
\DeclareMathSizes{9}{9}{5.5}{5}
\DeclareMathSizes{14.4}{14.4}{8.5}{6}
\setmainfont{Times New Roman}
\setsansfont{Arial}
\newfontfamily\coverfont{Arial}
\newfontfamily\calibrifont{Carlito}
\usepackage{unicode-math}
\setmathfont{TeX Gyre Termes Math}
\setmathfont[range={up/{latin,Latin,num}}]{Times New Roman}
\setmathfont[range={it/{latin,Latin}}]{Times New Roman Italic}
\setmathfont[range={bfup/{latin,Latin,num}}]{Times New Roman Bold}
\setmathfont[range={"002B,"003D,"003C,"003E}]{Standard Symbols PS}
\newfontfamily\pf@sym{Standard Symbols PS}
\newfontfamily\pf@symi[FakeSlant=0.21]{Standard Symbols PS}
% MathType-like spacing (tighter than TeX around = + -)
\AtBeginDocument{\thickmuskip=2.5mu plus 1mu\medmuskip=1.5mu plus 1mu\thinmuskip=3mu
% Symbol integral and sigma (18 pt at 12 pt text, as MathType), slanted Symbol pi
 \renewcommand{\int}{\mathop{\vcenter{\hbox{\pf@sym\fontsize{\fpeval{1.5*\f@size}}{0pt}\selectfont\char"F2}}}\nolimits}%
 \renewcommand{\sum}{\mathop{\vcenter{\hbox{\pf@sym\fontsize{\fpeval{1.5*\f@size}}{0pt}\selectfont\char"E5}}}\limits}%
 \renewcommand{\pi}{\mathord{\text{\pf@symi\char"70}}}}
\pfrestoreat
% Times New Roman + MathType-style maths, as his HW docs (de-docs/common/prelim_fonts.tex)
\usepackage{needspace,longtable,multicol,enumitem,array}
\providecommand{\vv}[1]{\overrightarrow{#1}}\providecommand{\cvv}[3]{\begin{pmatrix}#1\\#2\\#3\end{pmatrix}}
\raggedbottom
\setlength{\parindent}{0pt}
\fancyhf{}\lhead{\small Vectors I: Basics Worksheet}\rhead{\small Compiled by Maximus Pung Jun}
\cfoot{\small\thepage}
\renewcommand{\headrulewidth}{0.4pt}
\newcounter{wq}[section]
\newcommand{\Q}[1]{\par\needspace{5\baselineskip}\medskip\stepcounter{wq}\noindent\textbf{\thewq}\quad #1\par}
\newcommand{\wsec}[1]{\par\needspace{8\baselineskip}\bigskip\stepcounter{section}\noindent{\Large\scshape \thesection\quad #1}\par\medskip}
\renewcommand{\thesection}{\arabic{section}}
\begin{document}
\thispagestyle{fancy}
{\LARGE\scshape Vectors I Basics Worksheet: Questions}\par\smallskip
{Compiled by Maximus Pung Jun}\par\smallskip
{\small\textcolor{black!60}{Basics of vectors: $\mathbf i$, $\mathbf j$, $\mathbf k$ notation, magnitude, unit vectors, parallel vectors, collinear points, the ratio theorem, the scalar product and the vector product.}}\par\medskip
\wsec{$\mathbf{i}$, $\mathbf{j}$, $\mathbf{k}$ Notation and Column Vectors}
\Q{Write $\mathbf a=3\mathbf i-2\mathbf j+5\mathbf k$ as a column vector, and write $\mathbf b=\col{4}{0}{-1}$ in terms of $\mathbf i$, $\mathbf j$ and $\mathbf k$.}
\Q{The points $A$ and $B$ have coordinates $(1,2,3)$ and $(4,6,1)$. Find $\overrightarrow{AB}$ in terms of $\mathbf i$, $\mathbf j$ and $\mathbf k$.}
\begin{center}\begin{tikzpicture}[x={(0.8cm,0cm)},y={(0.4cm,0.32cm)},z={(0cm,0.8cm)},>=Stealth,line cap=round]
\draw[ax](0,0,0)--(5.8,0,0) node[right]{\scriptsize$x$};
\draw[ax](0,0,0)--(0,7,0) node[above right]{\scriptsize$y$};
\draw[ax](0,0,0)--(0,0,3.8) node[above]{\scriptsize$z$};
\node[below left] at (0,0,0){\scriptsize$O$};
\draw[dsh](1,2,3)--(1,2,0)--(1,0,0);\draw[dsh](1,2,0)--(0,2,0);
\draw[dsh](4,6,1)--(4,6,0)--(4,0,0);\draw[dsh](4,6,0)--(0,6,0);
\draw[va](0,0,0)--(1,2,3);\draw[vb](0,0,0)--(4,6,1);\draw[vc](1,2,3)--(4,6,1);
\node[dot,label=above:{\scriptsize$A(1,2,3)$}] at (1,2,3){};
\node[dot,label=right:{\scriptsize$B(4,6,1)$}] at (4,6,1){};
\end{tikzpicture}\end{center}
\Q{Given $\mathbf p=2\mathbf i+\mathbf k$ and $\mathbf q=-\mathbf i+3\mathbf j$, find $\mathbf r$ such that $2\mathbf p-3\mathbf q+\mathbf r=\mathbf 0$. Give $\mathbf r$ in terms of $\mathbf i$, $\mathbf j$ and $\mathbf k$.}
\Q{For the following statements, state whether it is \textbf{true} or \textbf{false}.}
\par\smallskip{\setlength{\extrarowheight}{5pt}\begin{longtable}{|>{\centering\arraybackslash}p{0.9cm}|p{\dimexpr\linewidth-0.9cm-1.6cm-6\tabcolsep-4\arrayrulewidth\relax}|>{\centering\arraybackslash}p{1.6cm}|}
\hline &\textbf{Statement}&\textbf{T/F}\\\hline\endhead
(a)&$\mathbf i+\mathbf j+\mathbf k=\col{1}{1}{1}$\rule[-8pt]{0pt}{8pt}&\\\hline
(b)&The $\mathbf j$-component of $4\mathbf i-\mathbf j+2\mathbf k$ is $-1$\rule[-8pt]{0pt}{8pt}&\\\hline
(c)&The position vector of $P(1,-2,0)$ is $\mathbf i-2\mathbf j$\rule[-8pt]{0pt}{8pt}&\\\hline
(d)&$3\mathbf i+2\mathbf k=\col{3}{2}{0}$\rule[-8pt]{0pt}{8pt}&\\\hline
(e)&$\overrightarrow{AB}=\overrightarrow{OA}-\overrightarrow{OB}$\rule[-8pt]{0pt}{8pt}&\\\hline
\end{longtable}}
\clearpage\wsec{Magnitude}
\Q{Find $|2\mathbf i-3\mathbf j+6\mathbf k|$.}
\Q{The points $A$ and $B$ have coordinates $(1,2,-1)$ and $(3,-4,2)$. Find the distance $AB$.}
\Q{Given $\mathbf a=\col{t}{2}{t-1}$ and $|\mathbf a|=3$, find the possible values of $t$.}
\Q{For the following statements, state whether it is \textbf{true} or \textbf{false}.}
\par\smallskip{\setlength{\extrarowheight}{5pt}\begin{longtable}{|>{\centering\arraybackslash}p{0.9cm}|p{\dimexpr\linewidth-0.9cm-1.6cm-6\tabcolsep-4\arrayrulewidth\relax}|>{\centering\arraybackslash}p{1.6cm}|}
\hline &\textbf{Statement}&\textbf{T/F}\\\hline\endhead
(a)&$|-\mathbf a|=|\mathbf a|$ for every vector $\mathbf a$\rule[-8pt]{0pt}{8pt}&\\\hline
(b)&$|3\mathbf a|=3|\mathbf a|$\rule[-8pt]{0pt}{8pt}&\\\hline
(c)&$\left|\col{0}{-5}{0}\right|=-5$\rule[-8pt]{0pt}{8pt}&\\\hline
(d)&$\left|\col{1}{1}{1}\right|=1$\rule[-8pt]{0pt}{8pt}&\\\hline
(e)&$|\mathbf i+2\mathbf j-2\mathbf k|=1+2-2=1$\rule[-8pt]{0pt}{8pt}&\\\hline
(f)&$|\mathbf a+\mathbf b|=|\mathbf a|+|\mathbf b|$ for all vectors $\mathbf a$, $\mathbf b$\rule[-8pt]{0pt}{8pt}&\\\hline
(g)&$|\mathbf a+\mathbf b|\le|\mathbf a|+|\mathbf b|$ for all vectors $\mathbf a$, $\mathbf b$\rule[-8pt]{0pt}{8pt}&\\\hline
(h)&For all $\lambda\in\mathbb R$, $|\lambda\mathbf a|=\lambda|\mathbf a|$\rule[-8pt]{0pt}{8pt}&\\\hline
\end{longtable}}
\clearpage\wsec{Unit Vectors}
\Q{Find the unit vector in the direction of $\col{2}{-1}{2}$.}
\Q{Find the vector of magnitude $10$ in the direction of $\col{3}{0}{-4}$.}
\Q{The points $A$ and $B$ have coordinates $(1,0,2)$ and $(3,-6,5)$. Find the vector of magnitude $14$ in the direction \emph{opposite} to $\overrightarrow{AB}$.}
\Q{For the following statements, state whether it is \textbf{true} or \textbf{false}.}
\par\smallskip{\setlength{\extrarowheight}{5pt}\begin{longtable}{|>{\centering\arraybackslash}p{0.9cm}|p{\dimexpr\linewidth-0.9cm-1.6cm-6\tabcolsep-4\arrayrulewidth\relax}|>{\centering\arraybackslash}p{1.6cm}|}
\hline &\textbf{Statement}&\textbf{T/F}\\\hline\endhead
(a)&$\col{1}{1}{1}$ is a unit vector\rule[-8pt]{0pt}{8pt}&\\\hline
(b)&$\dfrac13\col{1}{2}{2}$ is a unit vector\rule[-8pt]{0pt}{8pt}&\\\hline
(c)&$\mathbf i+\mathbf j$ is a unit vector\rule[-8pt]{0pt}{8pt}&\\\hline
(d)&$|\hat{\mathbf a}|=1$ for every non-zero vector $\mathbf a$\rule[-8pt]{0pt}{8pt}&\\\hline
(e)&$\mathbf a$ and $\hat{\mathbf a}$ are parallel\rule[-8pt]{0pt}{8pt}&\\\hline
\end{longtable}}
\clearpage\wsec{Parallel Vectors ($\mathbf a=\lambda\mathbf b$)}
\Q{Show that $\mathbf a=\col{2}{-1}{3}$ and $\mathbf b=\col{-6}{3}{-9}$ are parallel.}
\Q{Find $p$ and $q$ such that $\col{4}{p}{-2}$ is parallel to $\col{-2}{3}{q}$.}
\Q{Find all vectors of magnitude $6$ that are parallel to $\mathbf a=\col{1}{2}{-2}$.}
\Q{For the following statements, state whether it is \textbf{true} or \textbf{false}.}
\par\smallskip{\setlength{\extrarowheight}{5pt}\begin{longtable}{|>{\centering\arraybackslash}p{0.9cm}|p{\dimexpr\linewidth-0.9cm-1.6cm-6\tabcolsep-4\arrayrulewidth\relax}|>{\centering\arraybackslash}p{1.6cm}|}
\hline &\textbf{Statement}&\textbf{T/F}\\\hline\endhead
(a)&If $\mathbf a\parallel\mathbf b$, then $\mathbf a=\mathbf b$\rule[-8pt]{0pt}{8pt}&\\\hline
(b)&$\col{1}{-2}{3}$ is parallel to $\col{-2}{4}{-6}$\rule[-8pt]{0pt}{8pt}&\\\hline
(c)&$\col{2}{4}{6}$ is parallel to $\col{1}{2}{4}$\rule[-8pt]{0pt}{8pt}&\\\hline
(d)&$\col{3}{0}{-6}$ is parallel to $\col{-1}{0}{2}$\rule[-8pt]{0pt}{8pt}&\\\hline
(e)&If $\mathbf a=-2\mathbf b$, then $\mathbf a$ and $\mathbf b$ are not parallel since they point in opposite directions\rule[-8pt]{0pt}{8pt}&\\\hline
\end{longtable}}
\clearpage\wsec{Collinear Points}
\Q{Show that the points $A(1,2,3)$, $B(3,5,4)$ and $C(7,11,6)$ are collinear.}
\Q{For the points in Question 1, find the ratio $AB:BC$.}
\Q{The points $P(2,-1,0)$, $Q(4,1,3)$ and $R(a,7,b)$ are collinear. Find $a$ and $b$.}
\Q{For the following statements, state whether it is \textbf{true} or \textbf{false}.}
\par\smallskip{\setlength{\extrarowheight}{5pt}\begin{longtable}{|>{\centering\arraybackslash}p{0.9cm}|p{\dimexpr\linewidth-0.9cm-1.6cm-6\tabcolsep-4\arrayrulewidth\relax}|>{\centering\arraybackslash}p{1.6cm}|}
\hline &\textbf{Statement}&\textbf{T/F}\\\hline\endhead
(a)&If $\overrightarrow{AB}\parallel\overrightarrow{CD}$, then $A$, $B$, $C$ and $D$ are collinear\rule[-8pt]{0pt}{8pt}&\\\hline
(b)&If $\overrightarrow{AB}=3\overrightarrow{AC}$, then $A$, $B$ and $C$ are collinear\rule[-8pt]{0pt}{8pt}&\\\hline
(c)&$A(0,0,0)$, $B(1,1,1)$ and $C(2,2,3)$ are collinear\rule[-8pt]{0pt}{8pt}&\\\hline
(d)&The points with position vectors $\mathbf a$, $2\mathbf a$ and $-\mathbf a$ are collinear ($\mathbf a\neq\mathbf 0$)\rule[-8pt]{0pt}{8pt}&\\\hline
\end{longtable}}
\clearpage\wsec{Ratio Theorem}
\begin{center}\begin{tikzpicture}[scale=0.85,>=Stealth]
\coordinate (O) at (0,0);\coordinate (A) at (0.6,2.4);\coordinate (B) at (5.2,1.6);
\coordinate (C) at ($(A)!0.4!(B)$);
\draw[thick,black!70](A)--(B);
\draw[va](O)--(A) node[midway,left]{\small$\mathbf a$};\draw[vb](O)--(B) node[midway,below]{\small$\mathbf b$};\draw[vc](O)--(C);
\node[dot,label=below left:$O$] at (O){};\node[dot,label=above left:$A$] at (A){};
\node[dot,label=right:$B$] at (B){};\node[dot,label=below:$C$] at (C){};
\draw[decorate,decoration={brace,amplitude=4pt,raise=3pt}](A)--(C) node[midway,above=8pt,sloped]{$\lambda$};
\draw[decorate,decoration={brace,amplitude=4pt,raise=3pt}](C)--(B) node[midway,above=8pt,sloped]{$\mu$};
\end{tikzpicture}\end{center}\noindent In the True/False statements, $C$ lies on $AB$ with $AC:CB=\lambda:\mu$, $\overrightarrow{OA}=\mathbf a$ and $\overrightarrow{OB}=\mathbf b$.
\Q{Let $\overrightarrow{OA}=\mathbf i+\mathbf j+\mathbf k$ and $\overrightarrow{OB}=\mathbf i+2\mathbf j+3\mathbf k$. Given that $C$ lies on $AB$, find the coordinates of $C$ if\begin{enumerate}\item[(a)] $AC:CB=1:2$,\item[(b)] $3AC=2CB$,\item[(c)] $2AB=5CB$,\item[(d)] Given instead that $C$ lies on $AB$ produced, and $3AC=7BC$.\end{enumerate}}
\Q{Find the position vector of the midpoint $M$ of $AB$, where $A$ and $B$ have coordinates $(3,-1,5)$ and $(-1,5,1)$.}
\Q{For the following statements, state whether it is \textbf{true} or \textbf{false}.}
\par\smallskip{\setlength{\extrarowheight}{5pt}\begin{longtable}{|>{\centering\arraybackslash}p{0.9cm}|p{\dimexpr\linewidth-0.9cm-1.6cm-6\tabcolsep-4\arrayrulewidth\relax}|>{\centering\arraybackslash}p{1.6cm}|}
\hline &\textbf{Statement}&\textbf{T/F}\\\hline\endhead
(a)&$\overrightarrow{OC}=\dfrac{\mu\mathbf a+\lambda\mathbf b}{\lambda+\mu}$\rule[-8pt]{0pt}{8pt}&\\\hline
(b)&$\overrightarrow{OC}=\dfrac{\lambda\mathbf a+\mu\mathbf b}{\lambda+\mu}$\rule[-8pt]{0pt}{8pt}&\\\hline
(c)&If $\lambda:\mu=1:1$, then $\overrightarrow{OC}=\dfrac12(\mathbf a+\mathbf b)$\rule[-8pt]{0pt}{8pt}&\\\hline
(d)&If $\lambda:\mu=1:3$, then $\overrightarrow{OC}=\dfrac{3\mathbf a+\mathbf b}{4}$\rule[-8pt]{0pt}{8pt}&\\\hline
(e)&If $\lambda:\mu=2:1$, then $\overrightarrow{OC}=\dfrac{2\mathbf a+\mathbf b}{3}$\rule[-8pt]{0pt}{8pt}&\\\hline
\end{longtable}}
\clearpage\wsec{Scalar Product}
\begingroup\setlist{itemsep=1pt,topsep=2pt,parsep=0pt}\setlength{\abovedisplayskip}{3pt}\setlength{\belowdisplayskip}{3pt}
\Q{\textbf{Properties of scalar product}\par\smallskip It is given that $\mathbf a=\cvv{1}{2}{3}$, $\mathbf b=\cvv{-3}{2}{7}$, $\mathbf c=\cvv{4}{-5}{-1}$. Find the following\par\smallskip{\setlength{\tabcolsep}{0pt}\noindent\begin{tabular}{@{\hspace{1.2em}}*{4}{p{0.225\linewidth}}}(a)\ $\mathbf a\cdot\mathbf b$ & (b)\ $\mathbf b\cdot\mathbf a$ & (c)\ $\mathbf a\cdot(\mathbf b+\mathbf c)$ & (d)\ $\mathbf a\cdot\mathbf b+\mathbf a\cdot\mathbf c$\\[4pt](e)\ $\mathbf a\cdot(2\mathbf c)$ & (f)\ $(2\mathbf a)\cdot\mathbf c$ & (g)\ $\mathbf a\cdot\mathbf a$ &\end{tabular}}\par}
\Q{\textbf{Scalar product formula: $\mathbf a\cdot\mathbf b=|\mathbf a||\mathbf b|\cos\theta$}\par\smallskip \begin{enumerate}[label=(\alph*)]
\item Points $O$, $A$, $B$ are such that $\vv{OA}=\mathbf i+4\mathbf j-3\mathbf k$ and $\vv{OB}=5\mathbf i-\mathbf j$
\begin{enumerate}[label=(\roman*)]
\item Find the size of angle $AOB$
\item Find the size of angle $OAB$.
\end{enumerate}
\item It is given that $\mathbf c=\cvv{2}{3}{6}$, $\mathbf d=\cvv{1}{-2}{3}$ and $\mathbf e=\cvv{3}{2}{-2}$. Find the angle between the vectors $\mathbf c$ and
\begin{enumerate}[label=(\roman*)]
\item the vector $\mathbf d$
\item the vector $\mathbf e$
\item the positive $y$-axis
\end{enumerate}
\item It is given that two vectors $\mathbf a$ and $\mathbf b$ satisfy the following relationship:
\[
|\mathbf a||\mathbf b|=6.
\]
Find the value of $\mathbf a\cdot\mathbf b$ for each of the following scenarios.
\begin{enumerate}[label=(\roman*)]
\item The angle between $\mathbf a$ and $\mathbf b$ is $60^\circ$.
\item The angle between $\mathbf a$ and $\mathbf b$ is $120^\circ$.
\item Vectors $\mathbf a$ and $\mathbf b$ are perpendicular.
\end{enumerate}
\end{enumerate}}
\Q{\textbf{Perpendicularity of vectors}\par\smallskip \begin{enumerate}[label=(\alph*)]
\item Referred to the origin $O$, the position vectors of the points $A$ and $B$ are $\mathbf i-\mathbf j+2\mathbf k$ and $2\mathbf i+4\mathbf j+\mathbf k$ respectively. Show that $OA$ is perpendicular to $OB$.
\item The position vectors $\mathbf a$ and $\mathbf b$ are given by
\[
\mathbf a=2p\mathbf i+3p\mathbf j+6p\mathbf k\quad\text{and}\quad\mathbf b=\mathbf i-2\mathbf j+2\mathbf k
\]
where $p>0$. It is given that $|\mathbf a|=|\mathbf b|$.
\begin{enumerate}[label=(\roman*)]
\item Find the exact value of $p$.
\item Show that the vectors $(\mathbf a+\mathbf b)$ and $(\mathbf a-\mathbf b)$ are perpendicular.
\end{enumerate}
\end{enumerate}}
\Q{\textbf{Length of projection}\par\smallskip \begin{enumerate}[label=(\alph*)]
\item Given that $\mathbf a=2\mathbf i+\mathbf j+2\mathbf k$ and $\mathbf b=-\mathbf i+2\mathbf k-\mathbf k$, find the length of projection and the projection vector of
\begin{enumerate}[label=(\roman*)]
\item $\mathbf a$ onto $\mathbf b$,
\item $\mathbf b$ onto $\mathbf a$.
\end{enumerate}
\item With reference to the origin O, the points $A$, $B$ and $C$ have position vectors given by
\[
\mathbf a=\mathbf i+2\mathbf j-\mathbf k,\ \mathbf b=3\mathbf i+\mathbf k\ \text{and}\ \mathbf c=2\mathbf i+2\mathbf j+3\mathbf k
\]
respectively. Find the length of projection of $\vv{BA}$ onto $\vv{BC}$.
\end{enumerate}}
\endgroup
\clearpage\wsec{Vector Product}
\begingroup\setlist{itemsep=1pt,topsep=2pt,parsep=0pt}\setlength{\abovedisplayskip}{3pt}\setlength{\belowdisplayskip}{3pt}
\Q{\textbf{Properties of vector product}\par\smallskip \begin{enumerate}[label=(\alph*)]
\item Find $\cvv{3}{-1}{4}\times\cvv{5}{7}{1}$
\item Given that $\mathbf a=\cvv{1}{0}{-1}$, $\mathbf b=\cvv{1}{1}{1}$ and $\mathbf c=\cvv{-1}{1}{-1}$, find
\par\smallskip{\setlength{\tabcolsep}{0pt}\noindent\begin{tabular}{@{\hspace{1.2em}}*{4}{p{0.225\linewidth}}}(i)\ $\mathbf a\times\mathbf b$; & (ii)\ $\mathbf b\times\mathbf a$; & (iii)\ $\mathbf a\times(\mathbf b+\mathbf c)$; & (iv)\ $\mathbf a\times\mathbf b+\mathbf a\times\mathbf c$;\\[4pt](v)\ $\mathbf a\times(2\mathbf b)$; & (vi)\ $(2\mathbf a)\times\mathbf b$; & (vii)\ $\mathbf a\times\mathbf a$; &\end{tabular}}\par
\end{enumerate}}
\Q{\textbf{Parallel vectors}\par\smallskip \begin{enumerate}[label=(\alph*)]
\item Three non-zero vectors $\mathbf a$, $\mathbf b$ and $\mathbf c$ satisfy $\mathbf a\times3\mathbf b=2\mathbf a\times\mathbf c$. Show that $\mathbf a$ and $3\mathbf b-2\mathbf c$ are parallel.
\item Find the scalar product and vector product of the following vectors, stating whether they are parallel, perpendicular, or neither.
\begin{enumerate}[label=(\roman*)]
\item $\cvv{1}{2}{3}$ and $\cvv{4}{5}{6}$
\item $\mathbf a=\mathbf i+2\mathbf j-3\mathbf k$, $\mathbf b=3\mathbf i+\mathbf k$
\item $\cvv{1}{2}{-1}$ and $\cvv{-4}{-8}{4}$
\end{enumerate}
\end{enumerate}}
\Q{\textbf{Vector product formula}\par\smallskip Using the Vector Product, find the acute angles between the following vectors:\par\smallskip{\setlength{\tabcolsep}{0pt}\noindent\begin{tabular}{@{\hspace{1.2em}}*{3}{p{0.300\linewidth}}}(a)\ $\cvv{2}{1}{3}$ and $\cvv{4}{-1}{1}$ & (b)\ $\cvv{-1}{2}{-1}$ and $\cvv{3}{-1}{5}$ & (c)\ $\cvv{4}{3}{5}$ and $\cvv{2}{-5}{1}$\end{tabular}}\par}
\Q{\textbf{Perpendicular distance}\par\smallskip \begin{enumerate}[label=(\alph*)]
\item Relative to an origin $O$, two point $A$ and $B$ have position vectors given by $\mathbf a=3\mathbf i-\mathbf j+2\mathbf k$ and $\mathbf b=-2\mathbf i+3\mathbf j+2\mathbf k$. Find the following:
\begin{enumerate}[label=(\roman*)]
\item The shortest distance from $A$ to $OB$
\item The shortest distance from $B$ to $OA$
\end{enumerate}
\item Relative to the origin $O$, two points $A$ and $B$ have position vectors given by $\mathbf a=\mathbf i+\mathbf j+\mathbf k$ and $\mathbf b=\mathbf i+2\mathbf j+3\mathbf k$ respectively.
\begin{enumerate}[label=(\roman*)]
\item Find the length of projection of $\mathbf b$ onto $\mathbf a$.
\item Hence, or otherwise, find the shortest distance from $B$ to $OA$.
\end{enumerate}
\item Relative to the origin $O$, the points $A$, $B$ and $C$ have position vectors given by $\mathbf a=\mathbf i+2\mathbf j+3\mathbf k$, $\mathbf b=\mathbf i-\mathbf j+\mathbf k$ and $\mathbf c=3\mathbf i-\mathbf j+\mathbf k$. Find the shortest distance from $C$ to the line joining $A$ and $B$.
\end{enumerate}}
\Q{\textbf{Areas of triangles \& parallelograms}\par\smallskip \begin{enumerate}[label=(\alph*)]
\item Relative to the origin $O$, the points $A$ and $B$ have position vectors $\vv{OA}=\cvv{1}{0}{3}$, $\vv{OB}=\cvv{-2}{3}{2}$. Find the area of triangle $OAB$.
\item Relative to the origin $O$, the points $A$ and $B$ have position vectors $\vv{OA}=\cvv{0}{1}{1}$, $\vv{OB}=\cvv{1}{-2}{2}$. Find the area of triangle $OAB$.
\item Relative to an origin $O$, the position vectors of the vertices $A$, $B$ and $C$ of a triangle are
\[
\mathbf i+\mathbf j-\mathbf k,\ 2\mathbf i+\mathbf k\ \text{and}\ -\mathbf i+2\mathbf j-\mathbf k
\]
respectively. By considering the displacement vectors $\vv{AB}$ and $\vv{AC}$, find the exact area of triangle $ABC$.
\end{enumerate}}
\endgroup
\clearpage\wsec{Practice Problems}
\Q{Relative to the origin $O$, the position vectors of two points $A$ and $B$ are $\mathbf a$ and $\mathbf b$ respectively. The length of $\mathbf a$ is $4$ units and $\mathbf b$ is a unit vector. The angle between $\mathbf a$ and $\mathbf b$ is $\dfrac{\pi}{3}$ radians.\begin{enumerate}[label=\textbf{(\alph*)},leftmargin=2.2em,itemsep=3pt]\item Find the area of triangle $OAB$.\hfill[2]\end{enumerate}It is given that the point $M$ lies on the line segment $AB$ such that $AM:AB$ is $3:4$.\begin{enumerate}[label=\textbf{(\alph*)},leftmargin=2.2em,itemsep=3pt,start=2]\item Find the exact area of triangle $OAM$.\hfill[3]\end{enumerate}It is given that the point $C$ lies on $AB$ produced, such that $\overrightarrow{AC}=\lambda\overrightarrow{AB}$, where $\lambda>1$.\begin{enumerate}[label=\textbf{(\alph*)},leftmargin=2.2em,itemsep=3pt,start=3]\item Find the value of $\lambda$ such that $OC$ is perpendicular to $AB$.\hfill[4]\end{enumerate}}
\Q{The non-zero vectors $\mathbf a$ and $\mathbf b$ satisfy $|\mathbf a|=3$, $|\mathbf b|=4$ and $\mathbf a\cdot\mathbf b=6$.\begin{enumerate}[label=\textbf{(\alph*)},leftmargin=2.2em,itemsep=3pt]\item Find the angle between $\mathbf a$ and $\mathbf b$.\hfill[2]\item Find $|\mathbf a-\mathbf b|$.\hfill[3]\end{enumerate}}
\Q{Relative to the origin $O$, the points $A$, $B$ and $C$ have position vectors $\mathbf a$, $\mathbf b$ and $3\mathbf b$ respectively, where $\mathbf a$ and $\mathbf b$ are non-parallel.\begin{enumerate}[label=\textbf{(\alph*)},leftmargin=2.2em,itemsep=3pt]\item Find the area of triangle $OAC$ in terms of $|\mathbf a\times\mathbf b|$.\hfill[2]\item The point $D$ is the midpoint of $AC$. Find the area of triangle $OAD$ in terms of $|\mathbf a\times\mathbf b|$.\hfill[2]\end{enumerate}}
\Q{The non-zero vectors $\mathbf a$, $\mathbf b$ and $\mathbf c$ satisfy $3\mathbf a\times\mathbf b=2\mathbf c\times\mathbf b$.\begin{enumerate}[label=\textbf{(\alph*)},leftmargin=2.2em,itemsep=3pt]\item Show that $3\mathbf a-2\mathbf c=\lambda\mathbf b$ for some scalar $\lambda$.\hfill[2]\item It is now given that $|\mathbf a|=2$, $|\mathbf b|=3$ and the angle between $\mathbf a$ and $\mathbf b$ is $\dfrac{\pi}{6}$. Find $|3\mathbf a-2\mathbf b|$.\hfill[3]\end{enumerate}}
\Q{Relative to the origin $O$, the points $A$ and $B$ have position vectors $\mathbf a$ and $\mathbf b$ respectively. It is given that $|\mathbf a|=2$ and that $\mathbf b$ is a unit vector. The point $C$ divides the line segment $AB$ internally such that $AC:CB=3:2$. The point $D$ lies on $OB$ produced such that $\overrightarrow{OD}=3\overrightarrow{OB}$.\begin{enumerate}[label=\textbf{(\alph*)},leftmargin=2.2em,itemsep=3pt]\item Find $\overrightarrow{OC}$ in terms of $\mathbf a$ and $\mathbf b$.\hfill[2]\item Find $\overrightarrow{OD}$ in terms of $\mathbf b$.\hfill[1]\item Given that $CD$ is perpendicular to $AB$, find the value of $\mathbf a\cdot\mathbf b$.\hfill[3]\end{enumerate}}
\par\vspace{\baselineskip}\begin{center}\textbf{END OF WORKSHEET}\end{center}\end{document}
